% ch10-defense-demo — دليل المناقشة والعرض الحيّ
\section{دليل المناقشة والعرض: نتكلم عن إيه ونعرض إزاي}

الفصول السابقة شرحت \en{ScholarPath} من الداخل: الطبقات، و\en{CQRS}، والتوثيق،
والذكاء الاصطناعيّ، والفيديو، والمدفوعات، والأمان. هذا الفصل مختلف تمامًا: هو دليل
عمليّ يجهّزك لِلحظة الحقيقة — الوقوف أمام لجنة الدكاترة. لن نضيف هنا معلومة جديدة
عن النظام؛ بدل ذلك نجيب عن سؤالين لكل ميزة: «نتكلم فيها عن إيه؟» و«نعرضها حيّة
إزاي؟» — أي ماذا تفتح على الشاشة، وأين تضغط، وما الذي على اللجنة أن تراه بأمّ عينها.

لكل محور في هذا الفصل أربعة عناصر ثابتة: جملة تفتتح بها الكلام، نقاط القوّة التي
تتكلّم فيها، خطوات عرضٍ حيّ واقعية، وسؤال متوقّع من اللجنة مع إجابة قصيرة واثقة.
اقرأه مرّة قبل المناقشة، ثمّ استعمله كخريطة أثناءها.

% ====================================================================
\subsection{كيف تدير المناقشة؟}

مبدأ واحد يحكم الجلسة كلها، احفظه جيدًا: «لا تدّعِ شيئًا لا تستطيع إثباته على
الشاشة». اللجنة سمعت مئات العروض التي تصف ميزات لا وجود لها؛ ما يميّزك أنّ كل جملة
تقولها لها دليل — إمّا سطر كود تفتحه، أو زرّ تضغطه فيحدث الشيء أمامهم.

\begin{keyp}[نقطة مهمة: كل ادعاء له دليل]
رتّب كلامك على قاعدة «أقول ثمّ أُرِي». لا تقل «النظام يدعم الإشعارات اللحظيّة» ثمّ
تنتقل؛ قُلها ثمّ افتح نافذتين وابعث رسالة فيظهر الإشعار فورًا. الادعاء المصحوب
بعرضٍ حيّ يساوي عشرة ادعاءات بلا دليل، ويقطع الطريق على سؤال «طيّب وإثبات كده فين؟».
\end{keyp}

معظم عروضنا الحيّة تعتمد خدمات سحابية (\en{Azure OpenAI}, \en{ACS}, \en{Stripe}).
والسحابة قد تخذلك في أسوأ لحظة: النت يقطع، أو الخدمة تتأخّر. لهذا صُمِّم النظام منذ
البداية بمنطق «المنافذ والمحوّلات» (\en{ports/adapters})، وله لكل خدمة سحابية بديلٌ
محلّيّ (\en{Stub}) يشتغل \en{offline}.

\begin{note}[ملاحظة: خطة بديلة إن سقط النت]
لو انقطع الإنترنت أثناء ديمو سحابيّ، لا تتلعثم. اذكر بثقة أنّ النظام يُشغَّل بمنافذ
مجرّدة، وأنّ كل مزوّد خارجيّ (\en{Stripe}, \en{ACS}, \en{Azure OpenAI}) خلفه محوّل
\en{Stub} محلّيّ يعيد نتائج جاهزة بلا إنترنت — فتكمل عرض بقية الشاشات دون توقّف. هذا
نفسه نقطة قوّة معماريّة، لا اعتذار.
\end{note}

% ====================================================================
\subsection{خريطة عرضٍ في عشرين دقيقة}

قبل الدخول في التفاصيل، هذه خطّة زمنيّة تحكم إيقاع العرض كله. الفكرة أن توزّع
العشرين دقيقة على المحاور بحيث يبقى نصفها تقريبًا للعرض الحيّ لا للكلام. الجدول
يبيّن المدى الزمنيّ، والمحور، وما تفتحه على الشاشة في كل مرحلة:

\begin{center}
\begin{tabular}{C{2.2cm} L{5.2cm} L{7.0cm}}
\toprule
{\color{spPrimary}المدى الزمنيّ} & {\color{spPrimary}المحور} & {\color{spPrimary}تفتح إيه على الشاشة} \\
\midrule
0:00–3:00 & المقدمة والمشكلة & شرائح العرض: مشكلة تشتّت المنح والإرشاد، ثمّ الحلّ المقترَح \\
3:00–6:00 & المعماريّة & شجرة مجلّدات المشروع في المحرّر (\en{Domain/Application/Infrastructure/API}) \\
6:00–8:00 & التوثيق والتوكنات & المتصفّح مع \en{DevTools} على تبويبَي \en{Network} و\en{Application} \\
8:00–12:00 & الذكاء الاصطناعيّ (\en{RAG}) & صفحة إدارة قاعدة المعرفة $+$ نافذة الشات \\
12:00–15:00 & الفيديو (\en{ACS}) & نافذتان: حساب طالب وحساب مستشار في نفس الحجز \\
15:00–17:00 & الإشعارات والشات & نافذتان جنبًا لجنب لمستخدمَين مختلفَين \\
17:00–19:00 & المدفوعات (\en{Stripe}) & لوحة \en{Stripe} (وضع الاختبار) $+$ سجلّ الحجز \\
19:00–20:00 & الأمان & جدول \code{AuditLog} $+$ عمود مشفّر في القاعدة \\
\bottomrule
\end{tabular}
\end{center}

\begin{note}
راقب الساعة بطرف عينك. أكثر ما يفسد المناقشات هو الاستغراق في المقدمة حتى لا يبقى
وقت للعرض الحيّ — وهو بالضبط ما يريد الدكاترة رؤيته. لو ضاق الوقت، اختصر الكلام لا
الديمو.
\end{note}

% ====================================================================
\subsection{المحور الأول: المعماريّة (Clean Architecture و CQRS)}

افتتح بجملة: «قبل أي ميزة، نظامنا مبنيّ على معماريّة نظيفة تجعل كل جزء يعرف مكانه —
وهذا ما سمح لخمسة مطوّرين أن يعملوا متوازين بأقلّ تعارض ممكن».

نقاط القوّة التي تتكلّم فيها:
\begin{itemize}
\item الاعتماد يتّجه للداخل دائمًا: \en{Domain} في القلب لا يعرف شيئًا عن قاعدة
  البيانات ولا عن \en{ASP.NET}، والبنية التحتيّة هي التي تعتمد على القلب لا العكس.
\item الـ\en{controller} رفيع جدًّا: لا يحمل منطق عمل، فقط يمرّر الطلب.
\item كل حالة استخدام لها \en{handler} مستقلّ يملكه مطوّر واحد؛ فالتعديل في ميزة
  لا يلمس ملفّات ميزة أخرى، ممّا قلّل تعارضات الدمج (\en{merge conflicts}) بين
  الخمسة إلى أدنى حدّ.
\end{itemize}

\begin{idea}[اعرض كده: رِقّة الـ controller ثم الـ handler]
\begin{itemize}
\item افتح شجرة المجلّدات في المحرّر، ووسّع \code{server/src} لتُظهِر الطبقات
  الأربع جنبًا لجنب: \en{Domain}، \en{Application}، \en{Infrastructure}، \en{API}.
\item افتح أيّ \en{controller} وأشِر إلى أنّ الدالة سطر أو ثلاثة أسطر فقط، ثمّ افتح
  الـ\en{handler} المقابل لتُظهِر أن منطق العمل كلّه هناك، معزولًا وقابلًا للاختبار.
\end{itemize}
\end{idea}

المقتطف التالي يكفي لتوضيح رقّة الـ\en{controller}: إنّه يحوّل الطلب فقط إلى
\en{MediatR} ولا يقرّر شيئًا:

\begin{english}
\begin{lstlisting}[language={[Sharp]C}]
[HttpPost]
public async Task<IActionResult> Create(CreateBookingCommand cmd)
    => Ok(await sender.Send(cmd));   // the controller only forwards the request
\end{lstlisting}
\end{english}

\begin{qa}[سؤال متوقّع: CQRS مش overkill لمشروع تخرّج؟]
لا، وهذا ما يثبته حجم المشروع نفسه. عندنا نحو 211 متطلّبًا وخمسة مطوّرين؛ ومع هذا
الحجم كانت الحدود الواضحة (أمر لكل كتابة، استعلام لكل قراءة، \en{handler} مملوك
لمطوّر واحد) هي ما منع تعارضات الدمج وسمح بالعمل المتوازي. أضِف إلى ذلك أنّ نمط
الـ\en{pipeline} في \en{MediatR} أعطانا التحقّق والتسجيل والتدقيق «مجانًا» في مكان
واحد لكل الأوامر، بدل تكرارها في كل \en{controller}.
\end{qa}

% ====================================================================
\subsection{المحور الثاني: التوثيق والتوكنات (access مقابل refresh)}

هذا أكثر محور يُسأل عنه، فأتقِنه جيدًا. افتتح بجملة: «عند الدخول لا نصدر رمزًا
واحدًا بل رمزين، لكلٍّ دورٌ ونافذةُ خطرٍ مختلفة تمامًا».

نقاط القوّة:
\begin{itemize}
\item \en{Access JWT}: رمز حامل قصير العمر جدًّا — 15 دقيقة فقط — يُرفَق مع كل طلب،
  ويُفحَص محلّيًّا دون الرجوع للقاعدة (لذلك أبقيناه قصيرًا).
\item \en{Refresh token}: رمز عشوائيّ أطول عمرًا، وظيفته الوحيدة استبدال رمز الوصول
  المنتهي؛ يُدوَّر عند كل استخدام، ولا يُخزَّن خامًا بل مُجزَّأً (\en{SHA-256}) في
  الحقل \code{RefreshToken.TokenHash}.
\end{itemize}

\begin{idea}[اعرض كده: دورة حياة التوكن في DevTools]
\begin{itemize}
\item افتح \en{DevTools} على تبويب \en{Network}، سجّل الدخول، وأظهِر أن كل طلب يحمل
  ترويسة \code{Authorization: Bearer ...} فيها الـ\en{access JWT}.
\item انسخ الرمز إلى موقع فكّ \en{JWT} (أو أظهِر حمولته)، وأشِر إلى حقل انتهاء
  الصلاحيّة القصير.
\item انتظر انتهاء الصلاحيّة (أو أظهِر طلبًا يعيد \en{401})؛ عندها ينادي العميلُ
  \code{/api/auth/refresh} فيعود بزوجٍ جديد وتستمرّ الجلسة بلا إعادة إدخال كلمة المرور.
\item على تبويب قاعدة البيانات أظهِر أنّ عمود \code{RefreshToken.TokenHash} يحوي
  تجزئة (سلسلة \en{hex}) لا الرمز الخام.
\end{itemize}
\end{idea}

\begin{qa}[سؤال متوقّع: لو اتسرّب توكن يحصل إيه؟]
نفرّق بين الرمزين. لو تسرّب رمز الوصول فنافذته 15 دقيقة فقط ثمّ يصير عديم الفائدة.
ولو تسرّبت قاعدة البيانات فما فيها إلا تجزئات لا يمكن عكسها، فالتجزئة المسروقة لا
تُقبَل عند التحقّق. وفوق ذلك، رمز التحديث يُدوَّر عند كل استخدام مع كشف إعادة
الاستخدام: إن ظهر رمزٌ مُبطَل مرّة أخرى، نُلغي كامل عائلة رموز المستخدم فورًا كإجراء
ضدّ السرقة.
\end{qa}

\begin{qa}[سؤال متوقّع: ليه RS256 مش HS256؟]
لأنّ \en{RS256} يفصل التوقيع عن التحقّق: مفتاح خاصّ يوقّع، ومفتاح عامّ يتحقّق فقط.
الخدماتُ التي تتحقّق من الرمز تحتاج المفتاح العامّ وحده — ويمكن مشاركته بلا خطر لأنّه
لا يستطيع التوقيع. أمّا \en{HS256} فسرّه الواحد يوقّع ويتحقّق معًا، فمن يقرأه في أيّ
خدمة يستطيع تزوير رموز. \en{RS256} أيضًا يجعل تدوير المفاتيح عمليّة تشغيليّة سهلة.
\end{qa}

\begin{secnote}
الرسالة الأمنيّة هنا في جملة: نفصل «ما يُرسَل للعميل» عن «ما يُخزَّن عندنا». العميل
يأخذ الرمز الخام، والقاعدة تحفظ تجزئته فقط. هذا المبدأ نفسه نطبّقه على رمز إعادة
تعيين كلمة المرور، فتسريب القاعدة لا يمنح المهاجم أيّ رمزٍ قابلٍ للاستعمال.
\end{secnote}

% ====================================================================
\subsection{المحور الثالث: الذكاء الاصطناعيّ (RAG)}

افتتح بجملة: «مساعدنا الذكيّ لا يخترع الإجابات — إنّه يقرأ من قاعدة معرفة مبنيّة من
بيانات منصّتنا نفسها، ويستشهد بمصادره».

نقاط القوّة:
\begin{itemize}
\item خطّ \en{RAG} حقيقيّ: نبني قاعدة معرفة من بيانات المنصّة، نبحث فيها عن السياق
  الأقرب للسؤال، ثمّ نطلب من النموذج إجابةً مستندةً إلى ذلك السياق فقط.
\item الردّ يعرض «المصادر» أسفله، فيستطيع المستخدم التحقّق بنفسه.
\item هناك سقف تكلفة يوميّ يحمي الميزانيّة، كما أنّ ميزتَي التوصيات والأهليّة
  حسابيّتان محلّيًّا بلا أيّ نداء لِـ\en{GPT}.
\end{itemize}

\begin{idea}[اعرض كده: إعادة البناء ثم سؤال بمصادر]
\begin{itemize}
\item افتح صفحة إدارة قاعدة المعرفة، واضغط زرّ \en{Rebuild}، وأظهِر الأعداد العائدة:
  كم مستندًا \en{upserted} و\en{reembedded} و\en{removed} — فهذا يثبت أنّ البناء
  حيّ لا وهميّ.
\item انتقل إلى الشات واسأل سؤالًا عن منحة بعينها، وأظهِر الإجابة ثمّ كتلة «المصادر»
  أسفلها التي تشير إلى المستندات التي استندت إليها.
\item اذكر شفهيًّا سقفَ التكلفة اليوميّ، وأنّ التوصيات محسوبة محلّيًّا بلا \en{GPT}.
\end{itemize}
\end{idea}

\begin{qa}[سؤال متوقّع: المساعد بيهلوس (يخترع معلومات)؟]
احتمالُ الاختلاق مكبوح بطبقتين. أولًا، الـ\en{system prompt} يمنع النموذج صراحةً من
اختراع ما ليس في السياق، ويُلزمه أن يستند إلى كتلة سياقٍ مرقّمة تُمرَّر إليه. ثانيًا،
لو لم يوجد أيّ مستند تجاوز عتبة التشابه (نحو 0.6) نُعيد ردًّا عامًّا («لا تتوفّر لديّ
معلومة كافية») بدل التخمين. فالمساعد مقيَّد بالمصادر، لا يتكلّم من فراغ.
\end{qa}

\begin{qa}[سؤال متوقّع: ليه مش قاعدة بيانات متجهات (vector DB)؟]
لأنّ الحجم لا يبرّرها. عندنا نحو 692 مستندًا فقط، ومسحُها في الذاكرة لحساب التشابه
يستغرق أقلّ من 10 أجزاء من الألف من الثانية. إدخالُ قاعدة متجهات مستقلّة هنا يزيد
التعقيد والتكلفة التشغيليّة بلا مكسبٍ محسوس. لو كبرت القاعدة إلى مئات الآلاف من
المستندات، عندها يصبح القرار مختلفًا — والمعماريّة تسمح بالتبديل خلف منفذٍ واحد.
\end{qa}

% ====================================================================
\subsection{المحور الرابع: الجلسات المرئيّة (ACS)}

افتتح بجملة: «الاستشارة نفسها تحدث داخل المنصّة — مكالمة فيديو مسجَّلة بين الطالب
والمستشار، بلا الحاجة لأيّ أداة خارجيّة».

نقاط القوّة:
\begin{itemize}
\item مكالمة الفيديو تجري بين المتصفّحَين وخدمة \en{Azure Communication Services}؛
  خادمُنا لا يمرّر بايتات الصوت والصورة، بل ينسّق فقط.
\item لكل مشارك رمز وصول قصير العمر خاصّ به، والجلسة تُسجَّل تلقائيًّا.
\item نتتبّع الحضور بدقّة (لحظة دخول كلٍّ من الطالب والمستشار)، فنكتشف الغياب
  (\en{no-show}) آليًّا لا يدويًّا.
\end{itemize}

\begin{idea}[اعرض كده: جلسة حيّة من نافذتين]
\begin{itemize}
\item افتح حجزًا مؤكَّدًا، وادخل من نافذتين منفصلتين: واحدة بحساب الطالب وأخرى بحساب
  المستشار.
\item اضغط «انضمّ» في كلٍّ منهما، وأظهِر أنّ الاتصال قام، وأنّ كلّ طرفٍ يرى كاميرا
  الآخر ويسمعه.
\item أظهِر بدء التسجيل عند اتصال المكالمة، واذكر أنّ أوقات الدخول تُسجَّل وأنّ
  مهمّةً خلفيّة تعلّم أيّ حجزٍ لم يحضره طرفُه بـ\en{no-show}.
\end{itemize}
\end{idea}

\begin{qa}[سؤال متوقّع: ليه توكن لكل مشارك مش توكن واحد للغرفة؟]
لأنّ رمز الفيديو هو بيانات اعتماد: من يحمله يدخل المكالمة. الرمز المشترَك لا يمكن
إبطاله لطرفٍ دون الآخر ولا يميّز مَن دخل. أمّا رمزٌ قصير العمر لكل مشارك فيربط الدخول
بهويّة، ويجعل تتبّع الحضور دقيقًا، ويقلّص أثر أيّ تسريب إلى مستخدمٍ واحد ولدقائق
معدودة — فيصير الرمز المسروق عديم الفائدة بعد انتهاء الجلسة.
\end{qa}

% ====================================================================
\subsection{المحور الخامس: الإشعارات والشات (أقوى ديمو حيّ)}

هذا هو العرض الذي يُبهر اللجنة أكثر من غيره لأنّه لحظيّ ومرئيّ. افتتح بجملة: «كل شيء
في المنصّة يحدث في الزمن الحقيقيّ — الرسالة تصل، والجرس يرنّ، بلا تحديث صفحة».

\begin{idea}[اعرض كده: نافذتان جنبًا لجنب]
\begin{itemize}
\item افتح نافذتين لمستخدمَين مختلفَين جنبًا لجنب على الشاشة (يفضَّل نافذة عاديّة
  وأخرى خفيّة \en{incognito} لتفصل الجلستين).
\item ابعث رسالة شات من الأولى، وأظهِر ظهورها فورًا عند الثانية مع رنين الجرس
  (الإشعار) — دون أن تلمس زرّ التحديث.
\item فعّل «عدم الإزعاج» (\en{Do Not Disturb}) عند المستخدم الثاني، ثمّ ابعث إشعارًا
  آخر: أظهِر أنّه «يُخزَّن» في قائمة الإشعارات لكنّه «لم يقاطع» المستخدم — فالميزة
  تحترم تفضيله.
\end{itemize}
\end{idea}

\begin{qa}[سؤال متوقّع (أمان): إزاي تمنع مستخدمًا يتنصّت على محادثة مش بتاعته؟]
عبر بوّابة تمنع الوصول المباشر لمورد غير مملوك (\en{IDOR}) داخل \code{ChatHub}: قبل
أن ينضمّ أيّ عميلٍ إلى مجموعة محادثة، يتحقّق الـ\en{hub} أنّ صاحب الـ\en{JWT} هو
فعلًا أحد طرفَي تلك المحادثة، وإلّا يرفض الانضمام. أي أنّ الاشتراك نفسه محميّ، لا
مجرّد الدفع؛ فلا يستطيع مستخدمٌ مصادَقٌ الاشتراك في محادثة الآخرين والتنصّت عليها.
\end{qa}

\begin{secnote}
أضِف نقطةً تُظهِر عمق التفكير الأمنيّ: نصّ أيّ رسالة يُنقَّى من وسوم \en{HTML}
الخبيثة قبل تخزينه، لا بعد عرضه فقط. هذا دفاعٌ في العمق، لأنّ النصّ نفسه يسافر إلى
ثلاث وجهات — واجهة \en{React}، و\en{SignalR}، والبريد الإلكترونيّ — والتنقية عند
المصدر تحمي الوجهات الثلاث دفعةً واحدة.
\end{secnote}

% ====================================================================
\subsection{المحور السادس: المدفوعات (Stripe)}

افتتح بجملة: «فلوس الطالب تُحجَز ولا تُؤخَذ — لا تُخصَم فعليًّا إلا حين يقبل المستشار
الجلسة». هذا نموذج شبيه بالوسيط الأمين (\en{escrow}) مبنيّ على \en{Stripe}.

\begin{idea}[اعرض كده: تفويض ثم قبض]
\begin{itemize}
\item احجز جلسة كطالب، ثمّ افتح لوحة \en{Stripe} (وضع الاختبار) أو سجلّ الحجز في
  النظام.
\item أظهِر أنّ الـ\en{PaymentIntent} حالته \en{authorized} (تفويض): المبلغ محجوز
  على البطاقة لا مخصوم بعد.
\item بدّل إلى حساب المستشار واقبل الجلسة، ثمّ عُد لِلّوحة وأظهِر أنّ الحالة صارت
  \en{captured} — الآن فقط تحوّلت الفلوس فعليًّا.
\end{itemize}
\end{idea}

\begin{qa}[سؤال متوقّع: فين الفلوس بين الحجز والجلسة؟]
محجوزة على بطاقة الطالب عبر التفويض اليدويّ للقبض (\en{manual capture})، لا في
حسابنا ولا مخصومة منه. مصيرها يتحدّد بالحدث التالي: تُقبَض عند قبول المستشار،
تُفَكّ عند الرفض أو انتهاء المهلة، وتُسترَدّ (كاملة أو جزئيًّا حسب السياسة) عند
الغياب (\en{no-show}). فالطالب لا يخسر قرشًا مقابل خدمةٍ لم تُؤكَّد.
\end{qa}

\begin{secnote}
اذكر ثلاث ضمانات تُظهِر نضج التعامل مع المال: أوّلًا، كل عمليّة تحمل مفتاح
لا-تكرار (\en{idempotency}) فلا يُخصَم المبلغ مرّتين لو تكرّر الطلب. ثانيًا، نتحقّق من
توقيع كل \en{webhook} قادم من \en{Stripe} قبل تصديقه. ثالثًا، بيانات البطاقة لا تلمس
خوادمنا إطلاقًا — تُدخَل مباشرةً في عناصر \en{Stripe} الآمنة، فلا نخزّن رقم بطاقة قطّ.
\end{secnote}

% ====================================================================
\subsection{المحور السابع: الأمان (سريع ومؤثّر)}

الوقت هنا دقيقة واحدة، فاجعلها مركّزة على ثلاثة أدلّة ملموسة تُثبِت أنّ الأمان ليس
شعارًا. افتتح بجملة: «الأمان عندنا ليس ميزة مضافة، بل طبقات مدمجة يمكنك رؤيتها في
القاعدة نفسها».

\begin{idea}[اعرض كده: ثلاثة أدلّة أمنيّة]
\begin{itemize}
\item افتح جدول \code{AuditLog} وأظهِر صفًّا يسجّل عمليّة حسّاسة بقيمتها «قبل» و«بعد»
  ومَن نفّذها ومتى — فكل تغيير مهمّ متروكٌ له أثر.
\item أظهِر في القاعدة عمودًا مشفَّرًا: القيمة تظهر نصًّا مشفّرًا (\en{ciphertext})
  لا نصًّا مقروءًا، فحتى من يفتح القاعدة لا يقرأ البيانات الحسّاسة.
\item ارفع ملفًّا مصابًا (عيّنة اختبار) وأظهِر أنّ فحص \en{ClamAV} يرفضه — والسلوك
  «الفشل المغلق» (\en{fail-closed}): لو تعذّر الفحص لأيّ سبب، يُرفَض الملفّ لا يُقبَل.
\end{itemize}
\end{idea}

\begin{keyp}[نقطة مهمة: الأمان مرئيّ لا موصوف]
كل ادعاء أمنيّ في هذا المحور له أثرٌ ملموس على الشاشة — صفٌّ في سِجلّ، أو نصٌّ
مشفّر، أو ملفٌّ مرفوض. هذا أقوى بكثير من قول «النظام آمن»، ويُنهي المناقشة على انطباع
احترافيّ.
\end{keyp}

% ====================================================================
\subsection{نصائح ختاميّة للعرض}

قبل الجلسة بقليل، جهّز مسرح العرض حتى لا تضيّع دقائق ثمينة في الترتيب أمام اللجنة:

\begin{itemize}
\item سجّل الدخول مسبقًا بالحسابات التجريبيّة المزروعة في النظام قبل بدء المناقشة؛
  لا تخترع كلمات مرور على الهواء — استخدم الحسابات التجريبيّة المزروعة الجاهزة.
\item جهّز التبويبات مقدّمًا: تبويب لكل محور (المحرّر، صفحة قاعدة المعرفة، الشات،
  الحجز، لوحة \en{Stripe})، مرتّبة بترتيب خريطة العشرين دقيقة.
\item للديمو الحيّ ثنائيّ المستخدم (الفيديو، والشات)، رتّب النافذتين جنبًا لجنب على
  الشاشة قبل البدء حتى تظهر اللحظيّة بوضوح.
\item حضّر خطة بديلة لو انقطع النت: تذكّر أنّ لكل خدمة سحابية محوّل \en{Stub} محلّيّ،
  فتُكمِل عرض بقيّة الشاشات \en{offline} دون ارتباك.
\item لا تُطِل في المقدمة على حساب العرض الحيّ؛ اللجنة تريد أن ترى النظام يعمل.
\end{itemize}

\begin{note}
تدرّب على العرض كاملًا مرّة واحدة على الأقلّ بالساعة في يدك، من تسجيل الدخول حتى
الأمان. التدريب يكشف التبويب الناقص، والحساب غير المزروع، والخطوة التي تأخذ وقتًا
أطول ممّا تظنّ — قبل أن تكشفها اللجنة.
\end{note}

% ====================================================================
\begin{summary}[الخلاصة في سطور]
\begin{itemize}
\item أدِر المناقشة بقاعدة «أقول ثمّ أُرِي»: كل ادعاء يُدعَم بكودٍ أو عرضٍ حيّ، ولكل
  خدمة سحابية بديلٌ محلّيّ يعمل \en{offline} لو سقط النت.
\item وزّع العشرين دقيقة بخريطة واضحة، واترك نصفها تقريبًا للعرض الحيّ لا للكلام.
\item ثلاث رسائل يجب أن تبقى عالقة في ذهن اللجنة بعد خروجك:
  \begin{enumerate}
    \item تطويرٌ مدفوعٌ بالمواصفات (نحو 211 متطلّبًا) وكلّ شيءٍ مُدقَّق عبر
      \code{AuditLog} و\en{pipeline} موحَّد.
    \item ذكاءٌ اصطناعيّ إنتاجيّ بتكلفة حدّية قريبة من الصفر (توصيات محليّة، بحثٌ في
      الذاكرة) مع سقوطٍ آمن يمنع الاختلاق.
    \item بنية فيديو (\en{ACS}) وزمنٍ حقيقيّ (\en{SignalR}) متكاملة داخل المنصّة، مع
      أمانٍ مرئيّ لا موصوف (توكنات قصيرة، تشفير، فحص ملفّات، بوّابات \en{IDOR}).
  \end{enumerate}
\end{itemize}
\end{summary}
